\begin{figure}[p]
\caption{Examples of hint masking and truncation applied to \textbf{solution} reasoning traces. Left columns show the full hint; right columns show the result after applying each method at $h = 0.5$. Masking replaces randomly selected words with \texttt{[MASK]}; truncation retains only the first $h \times 100\%$ of words.}
\label{fig:solution-hint-examples}
\vspace{4pt}
\begin{tabular}{@{}p{0.48\textwidth}p{0.48\textwidth}@{}}
\textbf{\sffamily\small GPQA --- masking ($h = 0.5$)} \\[2pt]
\begin{tcolorbox}[
    colback=gray!3, colframe=blue!60!black, boxrule=0.5pt,
    title={\scriptsize\textbf{full hint}},
    fonttitle=\sffamily, coltitle=white, colbacktitle=blue!60!black,
    left=3pt, right=3pt, top=1pt, bottom=1pt
]
\begin{lstlisting}
To solve this problem, we need to identify what physical phenomenon causes the observed spectral feature and then calculate the corresponding comoving distance.

**Step 1: Identifying the Spectral Feature**

The observation shows a peak at 790 nm with a significant flux drop at shorter wavelengths. This is characteristic of the Lyman limit, which occurs when photons have enough energy to ionize neutral hydrogen from the ground state. In the rest [...]
\end{lstlisting}
\end{tcolorbox} & \begin{tcolorbox}[
    colback=gray!3, colframe=red!60!black, boxrule=0.5pt,
    title={\scriptsize\textbf{masking applied}},
    fonttitle=\sffamily, coltitle=white, colbacktitle=red!60!black,
    left=3pt, right=3pt, top=1pt, bottom=1pt
]
\begin{lstlisting}
[MASK] [MASK] [MASK] [MASK] we [MASK] to identify what [MASK] [MASK] [MASK] [MASK] [MASK] [MASK] feature and [MASK] calculate the corresponding [MASK] distance.

**Step 1: Identifying [MASK] [MASK] [MASK]

The observation [MASK] [MASK] peak [MASK] [MASK] nm [MASK] [MASK] significant flux [MASK] [MASK] shorter [MASK] [MASK] is characteristic of the Lyman limit, which occurs when photons [MASK] enough energy [MASK] ionize [MASK] [MASK] from [MASK] [MASK] [MASK] In [MASK] rest [MASK]
\end{lstlisting}
\end{tcolorbox} \\[6pt]
\textbf{\sffamily\small GPQA --- truncation ($h = 0.5$)} \\[2pt]
\begin{tcolorbox}[
    colback=gray!3, colframe=blue!60!black, boxrule=0.5pt,
    title={\scriptsize\textbf{full hint}},
    fonttitle=\sffamily, coltitle=white, colbacktitle=blue!60!black,
    left=3pt, right=3pt, top=1pt, bottom=1pt
]
\begin{lstlisting}
To solve this problem, we need to identify what physical phenomenon causes the observed spectral feature and then calculate the corresponding comoving distance.

**Step 1: Identifying the Spectral Feature**

The observation shows a peak at 790 nm with a significant flux drop at shorter wavelengths. This is characteristic of the Lyman limit, which occurs when photons have enough energy to ionize neutral hydrogen from the ground state. In the rest [...]
\end{lstlisting}
\end{tcolorbox} & \begin{tcolorbox}[
    colback=gray!3, colframe=red!60!black, boxrule=0.5pt,
    title={\scriptsize\textbf{truncation applied}},
    fonttitle=\sffamily, coltitle=white, colbacktitle=red!60!black,
    left=3pt, right=3pt, top=1pt, bottom=1pt
]
\begin{lstlisting}
To solve this problem, we need to identify what physical phenomenon causes the observed spectral feature and then calculate the corresponding comoving distance.

**Step 1: Identifying the Spectral Feature**

The observation shows a peak at
\end{lstlisting}
\end{tcolorbox} \\[6pt]
\textbf{\sffamily\small AIME --- masking ($h = 0.5$)} \\[2pt]
\begin{tcolorbox}[
    colback=gray!3, colframe=blue!60!black, boxrule=0.5pt,
    title={\scriptsize\textbf{full hint}},
    fonttitle=\sffamily, coltitle=white, colbacktitle=blue!60!black,
    left=3pt, right=3pt, top=1pt, bottom=1pt
]
\begin{lstlisting}
Let's denote the biking rate as $b$ km/h, the jogging rate as $j$ km/h, and the swimming rate as $s$ km/h. Since Ed and Sue have the same rates for each activity, we can use these variables for both.

**Setting up the equations:**

Ed bikes for 2 hours, jogs for 3 hours, and swims for 4 hours, covering 74 km:
$$2b + 3j + 4s = 74 \quad \text{...(1)}$$

Sue jogs for 2 hours, swims for 3 hours, and bikes for 4 hours, covering 91 km:
$$4b + 2j + 3s [...]
\end{lstlisting}
\end{tcolorbox} & \begin{tcolorbox}[
    colback=gray!3, colframe=red!60!black, boxrule=0.5pt,
    title={\scriptsize\textbf{masking applied}},
    fonttitle=\sffamily, coltitle=white, colbacktitle=red!60!black,
    left=3pt, right=3pt, top=1pt, bottom=1pt
]
\begin{lstlisting}
[MASK] denote [MASK] [MASK] [MASK] [MASK] [MASK] km/h, the [MASK] rate [MASK] $j$ [MASK] [MASK] the [MASK] [MASK] as $s$ [MASK] [MASK] [MASK] and [MASK] [MASK] the [MASK] [MASK] [MASK] each [MASK] we can use [MASK] variables for [MASK]

**Setting up the equations:**

Ed bikes for [MASK] hours, [MASK] for 3 hours, [MASK] [MASK] [MASK] 4 [MASK] [MASK] 74 km:
$$2b + [MASK] + [MASK] = 74 \quad \text{...(1)}$$

[MASK] jogs [MASK] 2 hours, swims [MASK] 3 hours, [MASK] bikes [MASK] [MASK] [MASK] [MASK] [MASK] [MASK]
[MASK] + [MASK] [MASK] 3s [MASK]
\end{lstlisting}
\end{tcolorbox} \\[6pt]
\textbf{\sffamily\small AIME --- truncation ($h = 0.5$)} \\[2pt]
\begin{tcolorbox}[
    colback=gray!3, colframe=blue!60!black, boxrule=0.5pt,
    title={\scriptsize\textbf{full hint}},
    fonttitle=\sffamily, coltitle=white, colbacktitle=blue!60!black,
    left=3pt, right=3pt, top=1pt, bottom=1pt
]
\begin{lstlisting}
Let's denote the biking rate as $b$ km/h, the jogging rate as $j$ km/h, and the swimming rate as $s$ km/h. Since Ed and Sue have the same rates for each activity, we can use these variables for both.

**Setting up the equations:**

Ed bikes for 2 hours, jogs for 3 hours, and swims for 4 hours, covering 74 km:
$$2b + 3j + 4s = 74 \quad \text{...(1)}$$

Sue jogs for 2 hours, swims for 3 hours, and bikes for 4 hours, covering 91 km:
$$4b + 2j + 3s [...]
\end{lstlisting}
\end{tcolorbox} & \begin{tcolorbox}[
    colback=gray!3, colframe=red!60!black, boxrule=0.5pt,
    title={\scriptsize\textbf{truncation applied}},
    fonttitle=\sffamily, coltitle=white, colbacktitle=red!60!black,
    left=3pt, right=3pt, top=1pt, bottom=1pt
]
\begin{lstlisting}
Let's denote the biking rate as $b$ km/h, the jogging rate as $j$ km/h, and the swimming rate as $s$ km/h. Since Ed and Sue have the same rates for each activity, we can use these variables for both.

**Setting up the equations:**

Ed bikes for
\end{lstlisting}
\end{tcolorbox} \\[6pt]
\end{tabular}
\end{figure}